% =========================================================================
% PLANTILLA MAESTRA: APUNTES TÉCNICOS PARA IMPRESIÓN LÁSER MONOCROMÁTICA
% =========================================================================
% Optimizada específicamente para impresoras láser monocromo con tóner (HP LaserJet P1006 / CB435A).
%
% Criterios de Ingeniería y Tipografía:
% 1. Cero tramado sucio: Todo el texto es 100% negro puro (#000000). Sin texto gris.
% 2. Ahorro de tóner y prevención de ondulación: Cero fondos negros plenos ni
%    tramados densos. Cajas con fondo blanco y barra lateral de alto contraste.
% 3. Ergonomía de lectura física: Formato a dos columnas (twocolumn) en A4 para
%    mantener entre 55 y 68 caracteres por línea (regla de oro de Bringhurst).
% 4. Encuadernación física y doble faz: Modo 'twoside' con 'bindingoffset' de 7 mm
%    para perforado de ganchos o anillado espiral sin perforar el texto.
% 5. Tipografía láser robusta: Bitstream Charter para el cuerpo (diseñada
%    específicamente para impresoras láser de 300-600 DPI, con serifas fuertes
%    que no se desvanecen en el tóner) y TeX Gyre Heros (Helvetica) para títulos.
% 6. Numeración exterior: Números de página en los bordes externos (LE,RO) para
%    localización rápida al hojear el apunte impreso con el pulgar.
% =========================================================================

\documentclass[9pt,a4paper,twoside]{extarticle}
\usepackage[utf8]{inputenc}
\usepackage[T1]{fontenc}

% --- Tipografía de Alto Rendimiento para Tóner ---
\usepackage{charter}                         % Serif robusta para lectura a 600 DPI
\usepackage{tgheros}                         % Sans-Serif nítida para títulos
\usepackage{courier}                         % Monospaced estándar legible para código
\usepackage[spanish,es-tabla,es-noquoting]{babel}
\usepackage{microtype}                       % Ajuste óptico de márgenes y espaciado
\usepackage{multicol}                        % Entorno de columnas profesional y flexible

% --- Geometría con Compensación de Encuadernación ---
\usepackage{geometry}
\geometry{
    top=1.3cm,
    bottom=1.3cm,
    inner=1.4cm,                             % Margen interno base
    outer=1.2cm,                             % Margen externo
    bindingoffset=7mm,                       % Espacio reservado para anillado o perforado
    headheight=11pt,
    headsep=4mm,
    footskip=6mm
}
\setlength{\columnsep}{6mm}                  % Separación clara entre columnas

\usepackage{amsmath,amssymb}
\usepackage{booktabs}
\usepackage{tabularx}
\usepackage{enumitem}
\usepackage{fancyhdr}
\usepackage{titlesec}
\usepackage{listings}
\usepackage[many]{tcolorbox}
\usepackage{hyperref}

% --- Espaciados y Párrafos ---
\setlength{\parindent}{0pt}
\setlength{\parskip}{2.5pt}

% --- Encabezados y Pies para Impresión Doble Faz (Duplex) ---
\pagestyle{fancy}
\fancyhf{}
\fancyhead[LE]{\fontsize{7.5pt}{9pt}\selectfont\sffamily\color{black!80} \nouppercase{\leftmark}}
\fancyhead[RO]{\fontsize{7.5pt}{9pt}\selectfont\sffamily\color{black!80} \nouppercase{\rightmark}}
\fancyfoot[LE,RO]{\fontsize{8pt}{9.5pt}\selectfont\sffamily\bfseries \thepage}
\fancyfoot[RE,LO]{\fontsize{7pt}{8.5pt}\selectfont\sffamily\color{black!60} Apuntes de Cátedra -- Uso Personal}
\renewcommand{\headrulewidth}{0.4pt}
\renewcommand{\footrulewidth}{0.3pt}

% Estilo para primera página (sin encabezado superior)
\fancypagestyle{firstpage}{
    \fancyhf{}
    \fancyfoot[C]{\fontsize{8pt}{9.5pt}\selectfont\sffamily\bfseries \thepage}
    \renewcommand{\headrulewidth}{0pt}
    \renewcommand{\footrulewidth}{0.3pt}
}

% --- Formato de Títulos (Sans-Serif y Líneas Finas en Negro Puro) ---
\titleformat{\section}
  {\fontsize{10pt}{12pt}\bfseries\sffamily\color{black}}
  {\thesection.}{0.4em}{}
  [\vspace{1pt}\titlerule]

\titlespacing*{\section}{0pt}{5pt}{2.5pt}

\titleformat{\subsection}
  {\fontsize{8.8pt}{10.8pt}\bfseries\sffamily\color{black!90}}
  {\thesubsection}{0.4em}{}

\titlespacing*{\subsection}{0pt}{3.5pt}{1.5pt}

\titleformat{\subsubsection}
  {\fontsize{8.2pt}{9.8pt}\bfseries\sffamily\color{black!80}}
  {\thesubsubsection}{0.4em}{}

\titlespacing*{\subsubsection}{0pt}{2.5pt}{1pt}

% --- Listas Compactas ---
\setlist[itemize]{leftmargin=*,noitemsep,topsep=1.5pt,label=\tiny$\blacksquare$}
\setlist[enumerate]{leftmargin=*,noitemsep,topsep=1.5pt}

% --- Entornos TColorBox Monocromáticos de Alto Contraste ---
% 1. Caja de Concepto / Teorema (Borde izquierdo grueso, fondo blanco puro)
\newtcolorbox{laserbox}[1][]{
    enhanced,
    breakable,
    colback=white,
    colframe=black,
    arc=0.8mm,
    boxrule=0.5pt,
    leftrule=3.2pt,
    title={\fontsize{8pt}{9.5pt}\selectfont\sffamily\bfseries #1},
    coltitle=black,
    attach boxed title to top left={yshift=-2mm, xshift=2mm},
    boxed title style={colback=white, boxrule=0.4pt, arc=0.5mm, top=1pt, bottom=1pt, left=3pt, right=3pt},
    left=5pt, right=5pt, top=4pt, bottom=3pt,
    before skip=3pt, after skip=3pt,
    fontupper=\fontsize{8pt}{9.8pt}\selectfont
}

% 2. Caja de Resumen / Examen (Trama sutil al 4% de gris, no genera ruido visual)
\newtcolorbox{summarybox}[1][]{
    enhanced,
    breakable,
    colback=black!4,
    colframe=black!80,
    arc=0.8mm,
    boxrule=0.6pt,
    title={\fontsize{8pt}{9.5pt}\selectfont\sffamily\bfseries $\blacktriangleright$\ #1},
    coltitle=black,
    left=5pt, right=5pt, top=3pt, bottom=3pt,
    before skip=3pt, after skip=3pt,
    fontupper=\fontsize{8pt}{9.8pt}\selectfont
}

% 3. Caja de Fórmula / Ecuación Destacada
\newtcolorbox{mathbox}[1][]{
    enhanced,
    colback=white,
    colframe=black,
    arc=0.5mm,
    boxrule=0.5pt,
    left=4pt, right=4pt, top=3pt, bottom=3pt,
    before skip=2pt, after skip=2pt
}

% --- Configuración de Código (Listings en Blanco y Negro) ---
\lstdefinestyle{laserstyle}{
    basicstyle=\ttfamily\fontsize{7.5pt}{9pt}\selectfont\color{black},
    keywordstyle=\bfseries\color{black},
    stringstyle=\itshape\color{black!90},
    commentstyle=\itshape\color{black!60},
    frame=single,
    framerule=0.4pt,
    rulecolor=\color{black!70},
    backgroundcolor=\color{white},
    breaklines=true,
    breakatwhitespace=false,
    tabsize=2,
    showstringspaces=false,
    aboveskip=2pt,
    belowskip=2pt,
    xleftmargin=3pt,
    xrightmargin=3pt
}

% Hipervínculos discretos para impresión (texto negro con línea de enlace punteada o invisible)
\hypersetup{
    colorlinks=false,
    pdfborder={0 0 0},
    pdfauthor={Hugo Joaquín Hidalgo},
    pdfcreator={Antigravity Laser Engine}
}

% =========================================================================
% INICIO DEL DOCUMENTO
% =========================================================================
\begin{document}
\thispagestyle{firstpage}

% --- Banner Superior de Ancho Completo ---
\begin{tcolorbox}[
    enhanced,
    colback=white,
    colframe=black,
    arc=1mm,
    boxrule=0.8pt,
    left=8pt, right=8pt, top=5pt, bottom=5pt,
    before skip=0pt, after skip=6pt
]
    \begin{minipage}[c]{0.74\textwidth}
        {\fontsize{12pt}{14pt}\selectfont\bfseries\sffamily [TÍTULO DEL APUNTE O MATERIA]}\\[2pt]
        {\fontsize{8.5pt}{10.5pt}\selectfont\sffamily [Subtítulo o Compendio de Unidades Temáticas]}
    \end{minipage}
    \hfill
    \begin{minipage}[c]{0.24\textwidth}
        \raggedleft\fontsize{7.5pt}{9pt}\selectfont\sffamily
        \textbf{UNJu} -- Facultad de Ingeniería\\
        \textbf{Edición:} Impresión Láser B/N\\
        \textbf{Fecha:} \today
    \end{minipage}
\end{tcolorbox}

% =========================================================================
% CONTENIDO A DOBLE COLUMNA
% =========================================================================
\begin{multicols*}{2}

\section{Introducción al Módulo}
Texto introductorio maquetado en dos columnas con tipografía Charter. Observe la longitud equilibrada de las líneas que facilita el escaneo visual rápido sin forzar la trayectoria ocular.

\begin{laserbox}[Concepto Clave / Definición]
Aquí se describe una definición de examen o principio fundamental. La caja posee un fondo blanco puro que no gasta tóner ni genera manchas de tramado en la impresora láser, y destaca mediante una línea negra lateral sólida de 3.2pt.
\end{laserbox}

\subsection{Ecuaciones y Fórmulas Numéricas}
Las ecuaciones matemáticas renderizan con total precisión gracias al motor AMS-LaTeX:
\begin{mathbox}
$$S(p) = \frac{T(1)}{T(p)} = \frac{1}{(1-f) + \frac{f}{p}} \le \frac{1}{s}$$
\end{mathbox}
Donde $s = 1 - f$ representa la fracción estrictamente secuencial del algoritmo.

\subsection{Bloque de Comandos o Algoritmos}
El resaltado de sintaxis para impresión láser monocromática prescinde de colores y utiliza \textbf{negritas} para palabras clave e \textit{itálicas} para comentarios:

\begin{lstlisting}[style=laserstyle,language=bash]
# Sincronizacion de datos hacia el servidor
rsync -avz --delete /origen/ joaquin@192.168.20.200:/destino/
\end{lstlisting}

\section{Tablas y Comparativas de Cátedra}
Las tablas utilizan la directriz tipográfica de \textit{Tufte} y \textit{Booktabs}: líneas horizontales limpias con grosores jerárquicos y ausencia total de líneas verticales que ensucien la lectura.

\begin{center}
\renewcommand{\arraystretch}{1.15}
\fontsize{7.5pt}{9pt}\selectfont
\begin{tabularx}{\columnwidth}{l c X}
\toprule
\textbf{Elemento} & \textbf{Nivel} & \textbf{Impacto Operativo} \\
\midrule
Servidor Principal & Crítico & Caída total del sistema ERP. \\
Red LAN / WiFi & Alto & Afecta a terminales de cobro. \\
Backup en Disco & Medio & Permite recuperación DRP. \\
\bottomrule
\end{tabularx}
\end{center}

\begin{summarybox}[Puntos Clave para Examen Oral]
\begin{itemize}
    \item Primer punto crítico evaluado por el docente.
    \item Segundo punto referente a la justificación técnica.
    \item Tercer punto con foco en normativa o cálculo.
\end{itemize}
\end{summarybox}

\end{multicols*}

\end{document}
